\documentclass[12pt,a4paper,oneside,openany]{book}
\usepackage[utf8]{inputenc}
\usepackage[portuguese]{babel}

% Fonte mais próxima de Times New Roman
\usepackage{mathptmx}

% Alterar margens do documento
\usepackage[top=30mm,bottom=35mm,inner=35mm,outer=25mm]{geometry}
\renewcommand{\baselinestretch}{1.15}
\usepackage{setspace}
\usepackage{fancyhdr}
\newcommand{\helv}{%
\fontfamily{mathptmx}\fontseries{b}\fontsize{10}{10}\selectfont}
%\pagestyle{fancy}
%\renewcommand{\chaptermark}[1]{\markboth{\MakeUppercase{\thechapter. #1 }}{}}
%\renewcommand{\sectionmark}[1]{\markright{\thesection\ #1}}
%\fancyhf{}
%\fancyhead[RO]{\bfseries\rightmark}
%\fancyhead[LE]{\bfseries\leftmark}
\setlength{\headheight}{20pt}
\fancyfoot[C]{\thepage}
\renewcommand{\headrulewidth}{0pt}
\renewcommand{\footrulewidth}{0pt}
\addtolength{\headheight}{0.5pt}
\fancypagestyle{plain}{
  \fancyhead{}
  \renewcommand{\headrulewidth}{0pt}  
  \fancyhead[L]{\helv \textbf{TÍTULO DO PROJETO}}
  \fancyhead[R]{\helv \textbf{Ciências da Engenharia Aeronáutica}}
  \fancyfoot[L]{\helv \textbf{NOME DO ALUNO}}
  \fancyfoot[R]{\helv \textbf{\monthyeardate, Universidade Atlântica}}
}

%\SetSinglespace{1.5}

%Links integrados
\usepackage{hyperref}

% Package que permite incluir imagens
\usepackage{graphicx}

% Comentar para não incluir a lista de figuras e tabelas no índice
\usepackage[nottoc]{tocbibind}

% Para citar referências
%\usepackage[square,numbers]{natbib}
%\bibliographystyle{abbrvnat}
%\bibliography{./bib/references.bib}
\usepackage[backend=biber,style=numeric,maxcitenames=1,autocite=inline,natbib=true]{biblatex}
\addbibresource{./bib/tese.bib}

%Indentar primeiro parágrafo
\usepackage{indentfirst}

% Incluir até à subsecção no índice
\setcounter{tocdepth}{2}
\setcounter{secnumdepth}{2}

% Permite remover os espaços entre itens
\usepackage{enumitem}

% Caption para imagens
\usepackage[font=footnotesize]{caption}

% Referenciar anexos
\usepackage[toc,page]{appendix}

%Adicionar pdfs
\usepackage{pdfpages}

%Alterar nome de capitulo para a referencia
\usepackage{fmtcount}

\usepackage{url}

%Para expressões Matemáticas
\usepackage{amsmath}
\usepackage{bm}

%Bibliotecas
\usepackage{biblatex}
\usepackage{csquotes}

%\usepackage{filecontents}
\usepackage{float}


\usepackage{subcaption}
\usepackage{eurosym}

\usepackage{booktabs}

\usepackage[T1]{fontenc}
%\usepackage[scaled]{helvet}
%\renewcommand*\familydefault{\sfdefault}

\usepackage{ragged2e}



\providecommand{\keywords}[1]{\textbf{Keywords:} #1}
\providecommand{\keywordsP}[1]{\textbf{Palavras chave:} #1}

\captionsetup{%
  figurewithin=none,
  tablewithin=none,
}
\usepackage{chngcntr}
\counterwithout{equation}{chapter}

\newcommand{\MONTH}{%
  \ifcase\the\month
  \or Janeiro% 1
  \or Fevereiro% 2
  \or Março% 3
  \or Abril% 4
  \or Maio% 5
  \or Junho% 6
  \or Julho% 7
  \or Agosto% 8
  \or Setembro% 9
  \or Outubro% 10
  \or Novembro% 11
  \or Dezembro% 12
  \fi}
\makeatletter
\makeatletter
\newcommand{\monthyeardate}{\MONTH\ \the\year}
\makeatother

%Remover "Capítulo X"
\usepackage{titlesec}
%\titleformat{\section}[frame]
%{\normalfont}
%{\filright
%\footnotesize
%\enspace SECÇÃO \thesection\enspace}
%{3pt}
%{\large\bfseries\filcenter}

\renewcommand{\thechapter}{\arabic{chapter}}
\titleformat{\chapter}[display]
{\bfseries\Large}
{\filleft\MakeUppercase{\chaptertitlename} \Huge\thechapter}
{4ex}
{
\vspace{2ex}%
\filright}
[\vspace{2ex}%
\titlerule]

\begin{document}
\let\clearpage\relax
\justifying
%\onehalfspacing
\begin{sloppy}

%    \makeatletter
%    \renewcommand{\@makechapterhead}[1]{%
%    \vspace*{50 pt}%
%    {\setlength{\parindent}{0pt} \raggedright \normalfont
%    \bfseries\Huge
%    \ifnum \value{secnumdepth}>1 
%        \if@mainmatter\thechapter.\ \fi%
%    \fi
%    #1\par\nobreak\vspace{40 pt}}}
%    \makeatother

\frontmatter
%-----------------------------------------------------------------
%
% Capa (FRENTE)
%
%------------------------------------------------------------------

%Capa
\begin{titlepage}
\include{./chapters/capa.tex}
\end{titlepage}



%-----------------------------------------------------------------
%
% Capa (TRÁS)
%
%------------------------------------------------------------------
%\clearpage \thispagestyle{empty}\mbox{}\clearpage
\pagenumbering{roman}
%\newpage
%\thispagestyle{plain}
%\mbox{}

%-----------------------------------------------------------------
%
% Rosto
%
%------------------------------------------------------------------

\include{chapters/rosto}

%-----------------------------------------------------------------
%
% Ideologia
%
%------------------------------------------------------------------
\newpage
\thispagestyle{plain}
\topskip0pt
\vspace*{\fill}
O autor é o único responsável pelas ideias expressas neste relatório.
\vspace*{\fill}

%-----------------------------------------------------------------
%
% AGRADECIMENTOS
%
%------------------------------------------------------------------
\newpage
\thispagestyle{plain}
\chapter*{Agradecimentos}
Obrigado.

%-----------------------------------------------------------------
%
% RESUMO
%
%------------------------------------------------------------------
\newpage
\thispagestyle{plain}
\chapter*{Resumo}

\keywordsP{um, dois, três, quatro, cinco}

%-----------------------------------------------------------------
%
% ABSTRACT
%
%------------------------------------------------------------------
\newpage
\thispagestyle{plain}
\chapter*{Abstract}
Range optimized electric UAV.
\keywords{one, two, three, four, five}

%-----------------------------------------------------------------
%
% ÍNDICE
%
%------------------------------------------------------------------
\newpage
\thispagestyle{plain}
\renewcommand{\contentsname}{Índice}
\tableofcontents
\newpage
\thispagestyle{plain}
\listoffigures
\newpage
\thispagestyle{plain}
\listoftables
%\clearpage \thispagestyle{plain}\mbox{}\clearpage

\mainmatter
%\pagenumbering{arabic} %numeros em numeração árabe
\pagestyle{plain}
%-----------------------------------------------------------------
%
% INTRODUÇÃO
%
%------------------------------------------------------------------
\newpage
\chapter{Introdução\label{ch:intro}}

\include{./chapters/01_intro.tex}

%-----------------------------------------------------------------
%
% Desenvolvimento
%
%-----------------------------------------------------------------
\newpage
\chapter{Desenvolvimento \label{ch:desenvolvimento}}

\include{./chapters/02_desenvolvimento.tex}

%-----------------------------------------------------------------
%
% Conclusão
%
%-----------------------------------------------------------------
\newpage
\chapter{Conclusão \label{ch:conclusao}}

\include{./chapters/03_conclusao.tex}

%-----------------------------------------------------------------
%
% Referências Bibliográficas
% 
%-----------------------------------------------------------------
%\renewcommand{\bibname}{Bibliografia}
%\bibliographystyle{babplain}      % basic style, author-year citations
%\bibliography{./bib/references.bib}
\printbibliography

\end{sloppy}
\end{document}